\section{Threat Model and Security Analysis}
\label{sec:threat}
We assume an adversary that can run nodes, register as a validator (paying bond + PoW), send
arbitrary messages, and partition the network temporarily, but cannot break ed25519 signatures
or SHA-256, and does not control a majority of the bonded stake.

\textbf{Block injection / poisoning.} Every header field is signed and validated; each shard is
committed by hash. An outsider cannot insert a block: forged signatures, unauthorized producers,
tampered bodies (Merkle/body-hash mismatch), wrong height or parent, future timestamps, and
forged log roots are all rejected (\S\ref{sec:eval} exercises seven such variants). A lying shard
server is detected by the committed shard hash and skipped.

\textbf{Data withholding.} A producer could publish a header but withhold body shards. The
availability-first canonical rule treats a block that cannot be reconstructed to its commitment
as not-yet-valid; reconstruction succeeds as long as $K$ honest holders exist, which the
replication factor $r$ and HRW spread are sized to ensure. Succinct DA
sampling~\cite{hallandersen2025foundations} would strengthen this from an operational check to a
probabilistic proof and is future work.

\textbf{Availability analysis.} Quantify recoverability under random node loss. If each node is
independently unavailable with probability $f$ and a shard has $r$ holders, that shard is lost
only if all its holders are down, w.p.\ $f^r$, so it is available w.p.\ $p=1-f^r$. A block of $T$
shards is recoverable iff at least $K$ shards survive:
\[
\Pr[\text{recoverable}] \;=\; \sum_{i=K}^{T}\binom{T}{i}p^{\,i}(1-p)^{\,T-i}.
\]
Table~\ref{tab:avail} evaluates this for the saturated shape (K16\,M8, $r{=}3$). A shard is lost
only when all three independently chosen holders are simultaneously down ($f^3$), and the code
still tolerates $M{=}8$ such losses; the block-loss probability is therefore dominated by
$\binom{T}{M+1}(f^3)^{M+1}$ and is negligible even at a severe $f{=}0.3$. This is the quantitative
reason fragmentation does not trade away durability: replication ($r$) protects each shard, and
parity ($M$) protects the block.

\begin{table}[t]
\centering
\caption{Recoverability at K16\,M8, $r{=}3$, $T{=}24$. Shard loss $f^3$; block lost iff $>M$ shards lost.}
\label{tab:avail}
\small
\begin{tabular}{rrr}
\toprule
$f$ & shard-loss $f^3$ & $\Pr[\text{block lost}]$ (bound) \\
\midrule
0.05 & $1.3\times10^{-4}$ & $<10^{-30}$ \\
0.10 & $1.0\times10^{-3}$ & $<10^{-20}$ \\
0.20 & $8.0\times10^{-3}$ & $<10^{-11}$ \\
0.30 & $2.7\times10^{-2}$ & $<10^{-6}$ \\
\bottomrule
\end{tabular}
\end{table}

\textbf{Correlated and colluding holders.} The durability argument above assumes node failures are
\emph{independent}; the $f^r$ term is exactly this assumption. It is the model's main soft spot.
If the $r$ holders of a shard are correlated (same operator, same data center, or a storage pool
under one administrator), losing or withholding that shard costs far less than $f^r$ suggests, and
a colluding set that holds all $r$ copies of enough shards can selectively withhold a block while
appearing online. HRW placement spreads holders by identity, not by operator, so it does not by
itself defeat an adversary who registers many identities behind one back end (the registration PoW
raises but does not eliminate that cost). Two defenses are partial today and strengthen with the
specified work: replication $r$ and parity $M$ raise the number of distinct holders that must
collude, and the storage-weighted election (built) plus future bond slashing for missed proofs make
holding-but-withholding economically self-defeating. We state this as a real limitation, not a solved problem.

\textbf{Sybil and long-range attacks.} Mass identities cost bond + registration PoW per
identity. Rewriting finalized history is impossible by construction (\S\ref{sec:fork}); a
heavier-but-late branch below the finality depth is refused. Above the finality depth, the
heaviest honest branch wins.

\textbf{Leader grinding.} Because leader election is
$\mathrm{HRW}(\textit{prevHash},h,t)$, an adversary might try to bias future elections by choosing
the content of a block it produces so that the resulting hash favors its own future rank. Two
factors limit this. First, the registration PoW seeds each identity's HRW key, so an attacker
cannot cheaply mint identities whose keys win specific slots. Second, influencing $\textit{prevHash}$
requires actually producing the parent block (already a won election) and only shifts the
\emph{next} slot probabilistically, not deterministically; grinding over many candidate blocks costs
hashing for a marginal bias. The storage-weighted election (\S\ref{sec:consensus}) further binds
odds to a quantity, proven storage, that an attacker cannot grind by reshaping a block: it must
actually answer storage-proof challenges to gain weight.

\textbf{Bond economics.} The bond must exceed the expected gain from misbehavior for slashing to
deter it. Registration requires a minimum bond, and the slash burns it while rewarding the reporter
a fraction, so an equivocating validator faces a certain loss against an uncertain gain. Setting the
minimum bond and slash fraction relative to block reward is a parameter-tuning question we leave to
deployment; the mechanism, not the specific values, is the contribution.

\textbf{Equivocation.} Double-signing is slashable with on-chain evidence, making it
economically irrational. \textbf{Eclipse/networking} attacks are out of scope for this paper and
mitigated only partially (peer diversity); we note them as a limitation. Table~\ref{tab:threat}
summarizes the model.

\begin{table}[t]
\centering
\caption{Threat model summary. \checkmark\,= mitigated in the prototype; $\circ$\,= partial.}
\label{tab:threat}
\small
\begin{tabular}{lll}
\toprule
Threat & Mechanism & Status \\
\midrule
Block injection      & signed header + per-shard hash      & \checkmark \\
Tampered body        & Merkle/body-hash commitment         & \checkmark \\
Lying shard server   & committed shard hash, skip+ban      & \checkmark \\
Data withholding     & availability-first canonical rule   & $\circ$ \\
Sybil identities     & bond + one-shot registration PoW    & \checkmark \\
Equivocation         & slashing with on-chain evidence     & \checkmark \\
Long-range rewrite   & depth-based finality ($F$ deep)     & \checkmark \\
Cross-network replay & chain-id bound into every tx        & \checkmark \\
Eclipse / routing    & peer diversity only                 & $\circ$ \\
\bottomrule
\end{tabular}
\end{table}
